\chapter{Complessità computazionale}
\label{chap:complessita}

L'esistenza di un algoritmo risolutivo, discussa nella
\cref{sec:fi-algoritmi}, non ne determina l'efficienza. La
\emph{complessità computazionale} studia le risorse necessarie
all'esecuzione in funzione della dimensione dell'ingresso: principalmente
tempo e memoria. L'analisi permette di confrontare algoritmi che
risolvono lo stesso problema e di prevederne il comportamento per
ingressi non ancora elaborati.\footnote{Materiale di riferimento:
M. Ceccarello e C. Pizzi, \emph{Cenni di complessità
computazionale}, dispense del corso di Fondamenti di Informatica.}

\section{Analisi sperimentale}
\label{sec:fi-complessita-sperimentale}

L'analisi sperimentale misura le prestazioni di un'implementazione
eseguendola su ingressi di dimensione diversa. Per esplorare un ampio
intervallo di dimensioni con poche misurazioni, si procede generalmente
raddoppiando la dimensione dell'ingresso a ogni prova. Se il tempo
impiegato per un ingresso di dimensione \(n\) è \(t(n)\), il rapporto
fra \(t(2n)\) e \(t(n)\) descrive l'effetto del raddoppio e permette
di stimare l'ordine di crescita.

Nel modello \(t(n)=an^b\), con \(a>0\), si ha
\[
  \frac{t(2n)}{t(n)}=2^b,
  \qquad
  b=\log_2\!\left(\frac{t(2n)}{t(n)}\right).
\]
Un rapporto prossimo a \(4\) suggerisce quindi una crescita quadratica;
un rapporto prossimo a \(2\) suggerisce una crescita lineare. La
stima non costituisce, da sola, una dimostrazione della complessità.

In Python, il modulo \texttt{time} offre \texttt{time.time()}, che
restituisce un istante espresso in secondi. Per misurare intervalli di
tempo è adatto \texttt{time.perf\_counter()}. La durata si ottiene sottraendo le letture effettuate prima e
dopo l'esecuzione. Per confrontare più dimensioni si possono registrare
i risultati in formato CSV, utilizzando i file di testo introdotti
nella \cref{sec:fi-python-file}. Un esempio completo di misurazione
sarà presentato dopo le due implementazioni di \texttt{reverse} (\cref{sec:fi-complessita-reverse-lineare}).

Il tempo misurato dipende anche dai valori dei dati, dall'hardware,
dall'interprete o compilatore, dal linguaggio, dall'implementazione e
dalle condizioni di esecuzione. Ripetere le misurazioni e considerarne valori aggregati
riduce l'influenza delle fluttuazioni, ma non elimina tali dipendenze.
Per studiare la crescita dei costi senza basarsi sulle misure
effettuate su una particolare macchina si ricorre all'analisi teorica.

\section{Analisi teorica e modello di costo}
\label{sec:fi-complessita-modello}

L'analisi teorica associa un costo alle operazioni eseguite, senza
misurarne direttamente la durata. Nel \emph{modello di costo unitario},
le operazioni primitive hanno costo costante: assegnamenti,
operazioni aritmetiche e logiche, confronti e accessi a una posizione
di una lista. Il costo dell'accesso per indice è giustificato
dall'organizzazione mediante array della \cref{sec:fi-strutture-array}.
Si tratta di un'astrazione, che presuppone operandi di
dimensione limitata e operazioni elementari eseguibili in tempo
costante.

Non tutte le istruzioni Python hanno costo unitario. Su una lista,
\texttt{len(lista)} richiede tempo costante, mentre la costruzione di
\texttt{[0] * n} richiede tempo proporzionale a \(n\). Anche il costo
di una funzione dipende dalle operazioni svolte al suo interno e non
può essere identificato con quello della sola chiamata.

Fissato il modello, si indica con \(T(x)\) il costo di esecuzione
sull'ingresso \(x\). Per ingressi di dimensione \(n\), il costo può
dipendere anche dal contenuto: si distinguono il caso migliore e il
caso peggiore e, quando è specificata una distribuzione degli ingressi,
il caso medio. Negli esempi seguenti, la dimensione è spesso
rappresentata dalla lunghezza \(n\) di una lista; quando necessario,
si considerano più parametri dimensionali.

Nel conteggio dei passi si sommano i costi dei blocchi eseguiti in
sequenza. Per un'istruzione condizionale si considera il ramo
effettivamente percorso; per un limite superiore si può maggiorare
il costo con quello del ramo più oneroso; per un limite inferiore
valido in entrambi i rami, con il costo di quello meno oneroso.
Il costo di un ciclo è la somma dei costi delle sue iterazioni,
che possono essere differenti. Per esempio, due cicli annidati,
entrambi con intervallo \texttt{range(n)}, eseguono il corpo interno
\(n^2\) volte; se il ciclo esterno usa \texttt{range(n)} con indice
\texttt{i} e quello interno usa \texttt{range(i)}, il numero diventa
\(\sum_{i=0}^{n-1}i=n(n-1)/2\).
Il modello può essere semplificato ulteriormente scegliendo come passi
base soltanto le operazioni rilevanti per il confronto in esame.

La \emph{complessità temporale} riguarda il numero di operazioni;
quella \emph{spaziale} riguarda la quantità di memoria richiesta. Le
due risorse devono essere valutate separatamente.

\section{Notazioni asintotiche}
\label{sec:fi-complessita-asintotica}

Per confrontare le funzioni di costo si studia il loro andamento per
\(n\) sufficientemente grande. Le notazioni asintotiche trascurano i
fattori moltiplicativi costanti e il comportamento su un numero finito
di dimensioni iniziali.

\begin{definition}[Notazioni \(O\), \(\Omega\) e \(\Theta\)]
\label{def:fi-notazioni-asintotiche}
Siano \(f,g\colon\N\longrightarrow[0,+\infty)\), con \(g(n)>0\)
per \(n\) sufficientemente grande. Si definiscono
\begin{align*}
  f\in O(g)
  &\iff \exists c>0\ \exists n_0\in\N:\
      f(n)\le c\,g(n)\quad\forall n\ge n_0,\\
  f\in\Omega(g)
  &\iff \exists c>0\ \exists n_0\in\N:\
      f(n)\ge c\,g(n)\quad\forall n\ge n_0,\\
  f\in\Theta(g)
  &\iff \exists c_1,c_2>0\ \exists n_0\in\N:\
      c_1g(n)\le f(n)\le c_2g(n)\quad\forall n\ge n_0.
\end{align*}
\end{definition}

La notazione \(O\) esprime un limite superiore asintotico, \(\Omega\)
un limite inferiore e \(\Theta\) entrambi. In particolare,
\[
  f\in\Theta(g)
  \iff f\in O(g)\ \text{e}\ f\in\Omega(g).
\]

La \cref{fig:fi-notazioni-asintotiche} mostra i tre casi.

\begin{figure}[htbp]
  \centering
  \begin{tikzpicture}[unipd/diagram,x=1cm,y=1cm]
  \begin{scope}[shift={(0,0)}]
    \draw[->] (0,0) -- (4.3,0) node[below] {$n$};
    \draw[->] (0,0) -- (0,3.8);
    \draw[dashed] (1.1,0) -- (1.1,3.4) node[below=2pt] {$n_0$};

    \draw[thick] plot[smooth] coordinates {(0.2,0.4) (1.0,1.0) (2.0,1.9) (3.6,3.1)};
    \node[above right] at (3.55,3.05) {\scriptsize $c\,g(n)$};

    \draw[thick] plot[smooth] coordinates {(0.2,0.2) (1.0,0.7) (2.0,1.2) (3.6,2.0)};
    \node[right] at (3.6,2.0) {\scriptsize $f(n)$};

    \node at (2.1,-0.7) {\scriptsize $f\in O(g)$};
  \end{scope}

  \begin{scope}[shift={(5.7,0)}]
    \draw[->] (0,0) -- (4.3,0) node[below] {$n$};
    \draw[->] (0,0) -- (0,3.8);
    \draw[dashed] (1.1,0) -- (1.1,3.4) node[below=2pt] {$n_0$};

    \draw[thick] plot[smooth] coordinates {(0.2,0.3) (1.0,0.8) (2.0,1.5) (3.6,2.5)};
    \node[right] at (3.6,2.5) {\scriptsize $f(n)$};

    \draw[thick] plot[smooth] coordinates {(0.2,0.2) (1.0,0.6) (2.0,1.1) (3.6,1.8)};
    \node[above right] at (3.55,1.8) {\scriptsize $c\,g(n)$};

    \node at (2.1,-0.7) {\scriptsize $f\in \Omega(g)$};
  \end{scope}

  \begin{scope}[shift={(11.4,0)}]
    \draw[->] (0,0) -- (4.3,0) node[below] {$n$};
    \draw[->] (0,0) -- (0,3.8);
    \draw[dashed] (1.1,0) -- (1.1,3.4) node[below=2pt] {$n_0$};

    \draw[thick] plot[smooth] coordinates {(0.2,0.3) (1.0,0.7) (2.0,1.3) (3.6,2.0)};
    \node[right] at (3.6,2.0) {\scriptsize $c_1g(n)$};

    \draw[thick] plot[smooth] coordinates {(0.2,0.5) (1.0,1.0) (2.0,1.8) (3.6,2.6)};
    \node[right] at (3.6,2.6) {\scriptsize $f(n)$};

    \draw[thick] plot[smooth] coordinates {(0.2,0.8) (1.0,1.4) (2.0,2.3) (3.6,3.2)};
    \node[right] at (3.6,3.2) {\scriptsize $c_2g(n)$};

    \node at (2.1,-0.7) {\scriptsize $f\in \Theta(g)$};
  \end{scope}
\end{tikzpicture}
  \caption{Interpretazione qualitativa delle notazioni asintotiche:
  in \(O(g)\) la funzione \(f\) è asintoticamente limitata superiormente
  da un multiplo di \(g\); in \(\Omega(g)\) è limitata inferiormente;
  in \(\Theta(g)\) è compresa fra due multipli positivi di \(g\).}
  \label{fig:fi-notazioni-asintotiche}
\end{figure}



\begin{example}[Analisi asintotica nei tre casi]
Si consideri un algoritmo di ordinamento che verifica inizialmente
se una lista di \(n\) elementi è già ordinata. In caso affermativo,
termina; altrimenti applica QuickSort, scegliendo sempre l'ultimo
elemento come pivot. L'algoritmo QuickSort è approfondito nella
\cref{sec:fi-ordinamento-rapido}.

Nel \emph{caso migliore}, la lista è già ordinata e il controllo
preliminare richiede un numero di confronti proporzionale a \(n\).
Pertanto,
\[
  T_{\mathrm{best}}(n)\in O(n),\qquad
  T_{\mathrm{best}}(n)\in\Omega(n),\qquad
  T_{\mathrm{best}}(n)\in\Theta(n).
\]

Nel \emph{caso medio}, assumendo una distribuzione uniforme delle
permutazioni di \(n\) valori distinti, QuickSort richiede mediamente
un numero di operazioni proporzionale a \(n\log n\). Il controllo
preliminare non modifica tale ordine di crescita, quindi
\[
  T_{\mathrm{avg}}(n)\in O(n\log n),\qquad
  T_{\mathrm{avg}}(n)\in\Omega(n\log n),\qquad
  T_{\mathrm{avg}}(n)\in\Theta(n\log n).
\]

Nel \emph{caso peggiore}, una lista ordinata in senso decrescente
determina partizioni fortemente sbilanciate. Il numero di operazioni
di QuickSort cresce quadraticamente:
\[
  T_{\mathrm{worst}}(n)\in O(n^2),\qquad
  T_{\mathrm{worst}}(n)\in\Omega(n^2),\qquad
  T_{\mathrm{worst}}(n)\in\Theta(n^2).
\]

I risultati si riassumono nella tabella seguente:
\begin{center}
\begin{tabular}{lccc}
  \toprule
  Caso & \(O\) & \(\Omega\) & \(\Theta\)\\
  \midrule
  Migliore & \(O(n)\) & \(\Omega(n)\) & \(\Theta(n)\)\\
  Medio & \(O(n\log n)\) & \(\Omega(n\log n)\)
        & \(\Theta(n\log n)\)\\
  Peggiore & \(O(n^2)\) & \(\Omega(n^2)\) & \(\Theta(n^2)\)\\
  \bottomrule
\end{tabular}
\end{center}

Le tre notazioni si applicano separatamente a ciascun caso.
Quando il limite superiore e quello inferiore hanno lo stesso
ordine di crescita, si ottiene una stima asintotica precisa
mediante \(\Theta\).
\end{example}

Nel linguaggio corrente, quando si dice soltanto che un algoritmo
è \(O(g(n))\), ci si riferisce spesso al caso peggiore; è però opportuno
specificarlo esplicitamente.

\section{Un algoritmo quadratico per l'inversione di una lista}
\label{sec:fi-complessita-reverse-quadratico}

Si consideri il problema di costruire una nuova lista contenente gli
elementi di una lista data in ordine inverso. Gli accessi e le
operazioni sulle liste sono descritti nella \cref{sec:fi-python-liste}.
La seguente implementazione inserisce ogni elemento all'inizio della
lista di uscita, spostando verso destra quelli già inseriti.

\begin{unipdcode}[language=Python]
def reverse_quadratico(lista):
    n = len(lista)                              # 1
    output = [0] * n                            # n + 1
    output_size = 0                             # 1
    for i in range(n):                          # n iterazioni
        for j in range(output_size, 0, -1):     # i iterazioni
            output[j] = output[j - 1]           # 3 (sottrazione, lettura, scrittura)
        output[0] = lista[i]                    # 2 (lettura, scrittura)
        output_size += 1                        # 1
    return output                               # escluso
\end{unipdcode}


All'inizio dell'iterazione di indice \(i\), i primi \(i\) elementi di
\texttt{output} contengono, in ordine inverso, i primi \(i\) elementi
della lista di ingresso. Il ciclo interno li sposta di una posizione;
l'assegnamento successivo inserisce \texttt{lista[i]} in posizione
zero. Dopo \(n\) iterazioni si ottiene quindi la lista invertita.

Per determinare il costo si adotta un conteggio semplificato, secondo
il modello della \cref{sec:fi-complessita-modello}. L'istruzione
\texttt{n = len(lista)} e gli assegnamenti semplici valgono un passo
ciascuno, mentre la costruzione di \texttt{[0] * n} costa \(n\) passi,
oltre all'assegnamento a \texttt{output}. Si trascurano i controlli
e gli aggiornamenti impliciti dei cicli e la restituzione finale.
L'inizializzazione costa pertanto \(n+3\) passi.

Nel ciclo interno l'istruzione
\texttt{output[j] = output[j - 1]} comporta una sottrazione,
un accesso in lettura e uno in scrittura: tre passi. Poiché
\texttt{output\_size} vale \(i\) all'inizio dell'iterazione esterna,
questa istruzione viene eseguita \(i\) volte. Le istruzioni successive
costano complessivamente altri tre passi: una lettura, una scrittura e
un incremento. Ne segue
\begin{align*}
  T(n)
    &=n+3+\sum_{i=0}^{n-1}(3i+3)\\
    &=n+3+3\frac{n(n-1)}2+3n\\
    &=\frac32n^2+\frac52n+3
     \in\Theta(n^2).
\end{align*}
L'espressione è esatta nel modello di conteggio adottato; non è il
numero effettivo di istruzioni macchina eseguite da Python. Il costo
quadratico dipende dagli spostamenti ripetuti, non dai valori degli
elementi. La lista di uscita occupa \(\Theta(n)\) spazio.

\section{Un algoritmo lineare per lo stesso problema}
\label{sec:fi-complessita-reverse-lineare}

Gli spostamenti non sono necessari se ciascun elemento viene scritto
direttamente nella propria posizione definitiva. Visitando gli indici
in ordine crescente, si copia l'elemento di indice \(n-1-i\) nella
posizione \(i\) della lista di uscita.

\begin{unipdcode}[language=Python]
def reverse_lineare(lista):
    n = len(lista)                               # 1
    output = [0] * n                             # n + 1
    for i in range(n):                           # n
        output[i] = lista[n - 1 - i]             # 4
    return output                                # escluso
\end{unipdcode}


All'inizio dell'iterazione di indice \(i\), i primi \(i\) elementi di
\texttt{output} contengono, in ordine inverso, gli ultimi \(i\) elementi
della lista di ingresso. Ogni iterazione colloca un nuovo elemento
nella posizione definitiva; dopo \(n\) iterazioni si ottiene quindi
la lista invertita.

L'inizializzazione costa \(n+2\) passi nel modello precedente. Ogni
iterazione esegue due sottrazioni, una lettura e una scrittura: quattro
passi, indipendenti da \(n\). Si ottiene quindi
\[
  T(n)=n+2+4n=5n+2\in\Theta(n).
\]
Entrambe le implementazioni producono lo stesso risultato e richiedono
\(\Theta(n)\) memoria per la lista di uscita, ma la seconda passa da
un costo temporale quadratico a uno lineare.

Per confrontarle sperimentalmente, il seguente programma registra in
\texttt{tempi.csv} la durata delle esecuzioni su liste di dimensione
crescente. Si assumono definite le due funzioni precedenti.

\begin{unipdcode}[language=Python]
import time

with open("tempi.csv", "w", encoding="utf-8") as file:
    file.write("algoritmo,n,secondi\n")
    for n in (500, 1000, 2000, 4000, 8000):
        dati = list(range(n))
        for funzione in (reverse_quadratico, reverse_lineare):
            inizio = time.perf_counter()
            funzione(dati)
            durata = time.perf_counter() - inizio
            file.write(f"{funzione.__name__},{n},{durata:.8f}\n")
\end{unipdcode}


Il file risultante contiene, per ogni misura, il nome dell'algoritmo,
la dimensione dell'ingresso e il tempo in secondi. Le misure
singole restano soggette alle fluttuazioni dell'esecuzione.
Per input sufficientemente grandi, il rapporto teorico fra i costi dei
due algoritmi cresce proporzionalmente a \(n\).

\section{Somme prefisse: evitare calcoli ripetuti}
\label{sec:fi-complessita-somme-prefisse}

Data una lista di numeri \(v=[v_0,\ldots,v_{n-1}]\), si vuole
costruire una lista \(s\) tale che
\[
  s_i=\sum_{j=0}^{i}v_j,\qquad 0\le i<n.
\]
Per esempio, da \([3,1,4,1,5]\) si ottiene
\([3,4,8,9,14]\). Una prima soluzione ricalcola per ogni indice
la somma dall'inizio; una seconda conserva la somma già raggiunta.

\noindent\hspace*{2mm}%
\begin{minipage}[t]{\dimexpr\linewidth/2-4mm\relax}
\vspace{0pt}
\textbf{Versione quadratica}
\begin{unipdcode}[language=Python,
  basicstyle=\ttfamily\scriptsize,
  numbers=none,
  xleftmargin=0pt,
  xrightmargin=0pt,
  framexleftmargin=0.3em]
def somme_prefisse_quadratiche(valori):
    n = len(valori)              # 1
    risultato = [0] * n          # n + 1
    for i in range(n):           # n iterazioni
        somma = 0                # 1
        for j in range(i + 1):   # i + 1 iterazioni
            somma += valori[j]   # 2 2 (lettura, aggiornamento)
        risultato[i] = somma     # 1
    return risultato             # escluso
\end{unipdcode}
\end{minipage}%
\hspace*{4mm}%
\begin{minipage}[t]{\dimexpr\linewidth/2-4mm\relax}
\vspace{0pt}
\textbf{Versione lineare}
\begin{unipdcode}[language=Python,
  basicstyle=\ttfamily\scriptsize,
  numbers=none,
  xleftmargin=0pt,
  xrightmargin=0pt,
  framexleftmargin=0.3em]
def somme_prefisse_lineari(valori):
    n = len(valori)              # 1
    risultato = [0] * n          # n + 1
    somma = 0                    # 1
    for i in range(n):           # n iterazioni
        somma += valori[i]       # 2 (lettura, aggiornamento)
        risultato[i] = somma     # 1
    return risultato             # escluso
\end{unipdcode}
\end{minipage}



Nella prima versione, a ogni iterazione di indice \(i\) viene
calcolata la somma dei primi \(i+1\) elementi della lista di ingresso,
ottenendo il valore \(s_i\). Il ciclo interno viene quindi eseguito
\(i+1\) volte e il numero complessivo degli addendi elaborati è
\[
  \sum_{i=0}^{n-1}(i+1)=\frac{n(n+1)}2.
\]
La complessità temporale è pertanto
\[
  T_{\mathrm{quadratico}}(n)\in\Theta(n^2).
\]

Nella seconda versione, all'inizio dell'iterazione di indice \(i\),
l'accumulatore \texttt{somma} contiene la somma dei primi \(i\)
elementi. L'aggiornamento lo porta a \(s_i\), che viene scritto
nella lista di uscita. Dopo \(n\) iterazioni si ottengono tutte
le somme prefisse, senza ricalcolare quelle precedenti.
Ogni iterazione richiede un numero costante di operazioni, quindi
\[
  T_{\mathrm{lineare}}(n)\in\Theta(n).
\]

Entrambe le implementazioni richiedono \(\Theta(n)\) spazio per
la lista di uscita e restituiscono una lista vuota quando
l'ingresso è vuoto.

\section{Massima differenza fra elementi}
\label{sec:fi-complessita-differenza}

Data una lista non vuota di numeri \(v=[v_0,\ldots,v_{n-1}]\), si
vuole determinare la massima differenza \(v_i-v_j\) fra due elementi,
senza vincoli sull'ordine degli indici, che possono anche coincidere.
Per esempio, da \([3,9,1,4]\) si ottiene \(9-1=8\).
Una prima soluzione esamina tutte le coppie di elementi; una seconda
sfrutta la relazione
\[
  \max_{0\le i,j<n}(v_i-v_j)
  =\max_{0\le i<n}v_i-\min_{0\le j<n}v_j.
\]
È quindi sufficiente individuare il minimo e il massimo della lista.

\noindent\hspace*{2mm}%
\begin{minipage}[t]{\dimexpr\linewidth/2-4mm\relax}
\vspace{0pt}
\textbf{Versione quadratica}
\begin{unipdcode}[language=Python,
  basicstyle=\ttfamily\scriptsize,
  numbers=none,
  xleftmargin=0pt,
  xrightmargin=0pt,
  framexleftmargin=0.3em]
def differenza_massima_quadratica(valori):
    if not valori:                             # 1
        raise ValueError("Lista vuota")
    massimo = 0                                # 1
    for i in range(len(valori)):               # n iterazioni
        for j in range(len(valori)):           # n iterazioni
            differenza = valori[i] - valori[j] # 4 (2 letture, sottrazione, assegnamento)
            if differenza > massimo:           # 1
                massimo = differenza           # 1 se vero
    return massimo                             # escluso
\end{unipdcode}
\end{minipage}%
\hspace*{4mm}%
\begin{minipage}[t]{\dimexpr\linewidth/2-4mm\relax}
\vspace{0pt}
\textbf{Versione lineare}
\begin{unipdcode}[language=Python,
  basicstyle=\ttfamily\scriptsize,
  numbers=none,
  xleftmargin=0pt,
  xrightmargin=0pt,
  framexleftmargin=0.3em]
def differenza_massima_lineare(valori):
    if not valori:                      # 1
        raise ValueError("Lista vuota")
    minimo = massimo = valori[0]        # 3 (lettura, 2 assegnamenti)
    for i in range(1, len(valori)):     # n - 1 iter.
        valore = valori[i]              # 2 (lettura, assegnamento)
        if valore < minimo:                    # 1
            minimo = valore             # 1 se vero
        if valore > massimo:            # 1
            massimo = valore            # 1 se vero
    return massimo - minimo             # 1
\end{unipdcode}
\end{minipage}


Nella prima versione, la variabile \texttt{massimo} conserva la
maggiore differenza fra quelle già esaminate e zero. Il valore
iniziale è valido poiché, scegliendo \(i=j\), si ottiene sempre
una differenza nulla. Dopo aver visitato tutte le coppie,
\texttt{massimo} contiene quindi il risultato richiesto.
I due cicli annidati esaminano \(n^2\) coppie e la complessità
temporale è
\[
  T_{\mathrm{quadratico}}(n)\in\Theta(n^2).
\]

Nella seconda versione, all'inizio dell'iterazione di indice \(i\),
le variabili \texttt{minimo} e \texttt{massimo} contengono
rispettivamente il minimo e il massimo dei primi \(i\) elementi.
I confronti con l'elemento corrente conservano tale proprietà;
dopo \(n-1\) iterazioni, la loro differenza fornisce il risultato
richiesto. Ogni iterazione richiede un numero costante di operazioni,
quindi
\[
  T_{\mathrm{lineare}}(n)\in\Theta(n).
\]

Entrambe le implementazioni richiedono \(\Theta(1)\) spazio ausiliario.
Le funzioni sollevano \texttt{ValueError} per una lista vuota e
restituiscono zero quando la lista contiene un solo elemento.

\clearpage

\section{Riconoscimento di elementi duplicati}
\label{sec:fi-complessita-duplicati}

Data una lista \(v=[v_0,\ldots,v_{n-1}]\), si vuole stabilire
se contiene elementi duplicati, ossia se esistono indici distinti
\(i,j\) tali che \(v_i=v_j\). Una prima soluzione confronta ogni
elemento con quelli successivi; una seconda ordina una copia della
lista e verifica la presenza di elementi consecutivi uguali.

\noindent\hspace*{2mm}%
\begin{minipage}[t]{\dimexpr\linewidth/2-4mm\relax}
\vspace{0pt}
\textbf{Confronto delle coppie}
\begin{unipdcode}[language=Python,
  basicstyle=\ttfamily\scriptsize,
  numbers=none,
  xleftmargin=0pt,
  xrightmargin=0pt,
  framexleftmargin=0.3em]
def contiene_duplicati_quadratico(valori):
    n = len(valori)                      # 1
    for i in range(n):                   # al più n iterazioni
        for j in range(i + 1, n):        # al più n-i-1 iterazioni
            if valori[i] == valori[j]:   # 3 (2 letture, confronto)
                return True              # escluso
    return False                         # escluso
\end{unipdcode}
\end{minipage}%
\hspace*{4mm}%
\begin{minipage}[t]{\dimexpr\linewidth/2-4mm\relax}
\vspace{0pt}
\textbf{Ordinamento preliminare}
\begin{unipdcode}[language=Python,
  basicstyle=\ttfamily\scriptsize,
  numbers=none,
  xleftmargin=0pt,
  xrightmargin=0pt,
  framexleftmargin=0.3em]
def contiene_duplicati_ordinando(valori):
    ordinati = sorted(valori)              # O(n log n)
    for i in range(len(ordinati) - 1):     # al più n-1 iterazioni
        if ordinati[i] == ordinati[i + 1]: # 4 (somma, 2 letture, confronto)
            return True                    # escluso
    return False                           # escluso
\end{unipdcode}
\end{minipage}

Nella prima versione, il ciclo interno confronta l'elemento di indice
\(i\) con tutti quelli di indice maggiore. Ogni coppia di indici
distinti viene quindi esaminata al più una volta e la funzione
restituisce \texttt{True} non appena individua due valori uguali.
Nel caso peggiore, quando tutti gli elementi sono distinti,
il numero di confronti è
\[
  \sum_{i=0}^{n-1}(n-i-1)
  =\frac{n(n-1)}2.
\]
La complessità temporale nel caso peggiore è pertanto
\[
  T_{\mathrm{quadratico}}(n)\in\Theta(n^2).
\]
Nel caso migliore, per \(n\ge2\), i primi due elementi coincidono
e la funzione termina dopo un solo confronto, in tempo \(\Theta(1)\).

Nella seconda versione, l'ordinamento colloca gli eventuali valori
duplicati in posizioni adiacenti. Il ciclo verifica ciascuna coppia
di elementi consecutivi e restituisce \texttt{True} quando trova
un'uguaglianza. Se la scansione termina senza individuarne, tutti
gli elementi sono distinti. L'ordinamento mediante \texttt{sorted()}
richiede \(O(n\log n)\) tempo nel caso peggiore, mentre la scansione
richiede \(O(n)\). La complessità temporale complessiva è quindi
\[
  T_{\mathrm{ordinamento}}(n)\in O(n\log n).
\]

La prima implementazione richiede \(\Theta(1)\) spazio ausiliario,
mentre la seconda richiede \(O(n)\) spazio aggiuntivo per la copia
ordinata. Entrambe restituiscono \texttt{False} per liste con meno
di due elementi. La seconda soluzione richiede inoltre che i valori
siano ordinabili secondo un criterio coerente con l'uguaglianza.

\clearpage

\section{Ricerca di una sottostringa e dimensioni multiple}
\label{sec:fi-complessita-sottostringa}

Date due stringhe \texttt{testo} e \texttt{modello}, si vuole
restituire il primo indice in cui il modello compare nel testo,
oppure \(-1\) in caso di assenza. Per esempio, cercare
\texttt{"cad"} in \texttt{"abracadabra"} restituisce \(4\).
Indicando con \(n\) e \(m\) le rispettive lunghezze, i possibili
indici iniziali sono \(0,\ldots,n-m\), purché \(m\le n\).
La ricerca diretta confronta i caratteri del modello con quelli
del testo per ogni posizione candidata.

\begin{unipdcode}[language=Python]
def trova_sottostringa(testo, modello):
    n = len(testo)                             # 1
    m = len(modello)                           # 1
    for i in range(n - m + 1):                 # al più n-m+1 
        coincidenze = 0                        # 1
        for j in range(m):                     # m 
            if testo[i + j] == modello[j]:     # 4
                coincidenze += 1               # 1 (t)
        if coincidenze == m:                   # 1
            return i                           # escluso
    return -1                                  # escluso
\end{unipdcode}


Per ogni posizione iniziale \(i\), la variabile
\texttt{coincidenze} conta i caratteri corrispondenti nelle due
stringhe. Se il conteggio raggiunge \(m\), tutti i caratteri
del modello coincidono e la funzione restituisce \(i\).
Poiché le posizioni vengono esaminate in ordine crescente,
l'indice restituito è quello della prima occorrenza.
Se nessuna posizione soddisfa la condizione, viene restituito \(-1\).

Per \(1\le m\le n\), nel caso peggiore vengono esaminate tutte
le \(n-m+1\) posizioni candidate, confrontando \(m\) caratteri
per ciascuna. La complessità temporale è pertanto
\[
  T_{\mathrm{peggiore}}(n,m)
  \in\Theta\bigl((n-m+1)m\bigr).
\]

Se \(m\) è costante, il costo è \(\Theta(n)\); se \(m=n-1\),
soltanto due posizioni sono candidate e il costo è ancora
\(\Theta(n)\). Quando invece \(m\) è circa \(n/2\), il costo
diventa \(\Theta(n^2)\).

Nel caso migliore, quando la sottostringa cercata compare all'inizio
del testo, viene esaminata una sola posizione candidata, confrontando
\(m\) caratteri. La complessità temporale è pertanto
\[
  T_{\mathrm{migliore}}(n,m)
  \in\Theta(m).
\]

L'implementazione richiede \(\Theta(1)\) spazio ausiliario.
Se il modello è più lungo del testo, restituisce \(-1\) senza
eseguire confronti fra caratteri; se è vuoto, restituisce
zero, considerando la stringa vuota presente all'inizio del testo.

L'algoritmo potrebbe interrompere il ciclo interno al primo
carattere diverso, evitando confronti inutili senza eliminare
il costo quadratico del caso peggiore. Algoritmi più sofisticati,
come Knuth--Morris--Pratt, possono ottenere un tempo
\(O(n+m)\) nel caso peggiore.

\clearpage

\section{Ordini di crescita e trattabilità}
\label{sec:fi-complessita-ordini}

Fra gli ordini di crescita più comuni si distinguono quello costante
\(\Theta(1)\), logaritmico \(\Theta(\log n)\), lineare \(\Theta(n)\),
lineare-logaritmico \(\Theta(n\log n)\), quadratico \(\Theta(n^2)\)
e, più in generale, polinomiale \(\Theta(n^k)\) per \(k\) intero
positivo fissato. Una crescita esponenziale come \(\Theta(2^n)\)
supera asintoticamente qualsiasi polinomio di grado fissato;
la funzione \(n^n\) cresce ancora più rapidamente.

La base del logaritmo è irrilevante nelle notazioni asintotiche:
per \(a,b>1\), la formula del cambiamento di base
\[
  \log_a n=\frac{\log_b n}{\log_b a}
\]
mostra che i logaritmi differiscono soltanto per un fattore
moltiplicativo costante. Pertanto
\(\log_a n\in\Theta(\log_b n)\).

La distinzione fra crescita polinomiale ed esponenziale è importante
per la trattabilità dei problemi, ma non determina da sola la
praticabilità di un'implementazione: incidono anche la dimensione
concreta dell'ingresso, i coefficienti e le risorse disponibili.
In particolare, un algoritmo polinomiale con esponente elevato può
essere impraticabile per dimensioni moderate.

\section{Complessità della ricorsione di Fibonacci}
\label{sec:fi-complessita-fibonacci}

Le due implementazioni di Fibonacci sono presentate nella
\cref{sec:fi-fibonacci}, dove si è stabilito che il numero complessivo
delle invocazioni ricorsive è \(C_n=2F_n-1\).
Nel modello a costo unitario, ogni invocazione compie un numero
costante di operazioni, escluse quelle delle chiamate interne.
La complessità temporale della versione ricorsiva è pertanto
\[
  T_{\mathrm{ric}}(n)\in\Theta(C_n)=\Theta(F_n).
\]
Dalla ricorrenza e dalla monotonia della successione seguono
\(F_{n+2}\ge2F_n\) e \(F_{n+1}\le2F_n\) per \(n\ge2\).
Queste disuguaglianze permettono di ricavare le stime
\[
  2^{\lfloor(n-1)/2\rfloor}\le F_n\le2^{n-2},
\]
che mostrano la crescita esponenziale del numero delle invocazioni.
Il costo dipende dal ricalcolo degli stessi sottoproblemi, visibile
nella \cref{fig:fi-fibonacci-albero}.

La versione iterativa, invece, calcola ogni termine una sola volta
ed esegue un numero di operazioni proporzionale a \(n\). La sua
complessità temporale è quindi
\[
  T_{\mathrm{iter}}(n)\in\Theta(n).
\]

La versione iterativa richiede \(\Theta(1)\) spazio ausiliario,
mentre quella ricorsiva raggiunge una profondità di pila
\(\Theta(n)\) e richiede quindi \(\Theta(n)\) spazio.
Come nel confronto del fattoriale della
\cref{sec:fi-pila-ricorsione}, queste stime considerano
operazioni a costo unitario e non includono il costo crescente
dell'aritmetica su interi arbitrariamente grandi.
